\chapter{Introduction}

\section{Background and Motivation}

Quick-commerce (q-commerce) platforms such as Blinkit, Zepto, and
Instamart have rewired Indian urban delivery by promising arrival
times measured in minutes rather than hours. The operating pattern is
a network of small ``dark stores''---also called depots or
forward-deployed micro-warehouses---each serving a 2--3~km radius
through a fleet of two-wheeler riders. Order batches are dispatched
every few minutes, with each rider typically carrying one to two
items per route to keep the delivery promise realistic.

The combination of strict service-level agreements (SLAs), high
dispatch frequency, narrow residential-street penetration, and a
vulnerable rider population (helmeted but exposed two-wheelers in
mixed traffic) creates a routing problem in which safety is not an
optional refinement. India's road-safety burden makes this concrete:
the Ministry of Road Transport and Highways reported $4{,}61{,}312$
road accidents and $1{,}68{,}491$ fatalities in 2022
\citep{morth2022}, with \citet{mohan2017} showing that fatalities are
disproportionately concentrated on higher-order roads despite their
relatively small share of total network length.

\section{Why the Classical VRPTW is Insufficient}

The classical Vehicle Routing Problem with Time Windows
\citep{solomon1987,desrochers1992,toth2002} captures none of the
q-commerce setting. It minimises a weighted sum of travel time,
distance, or vehicle count subject to capacity and time windows.
Crash risk, when modelled \citep{hoseinzadeh2020}, is typically
embedded as an additional weighted term in the objective using a
constant ``risk-per-edge'' coefficient that has no calibration
story. Several works in safety-aware routing penalise hazardous arcs
implicitly via inflated travel times, which conflates two distinct
phenomena (slow vs.\ dangerous) and gives the operator no explicit
knob.

Two structural problems follow. First, weighted-sum aggregation
collapses incompatible quantities (minutes, expected crash counts,
counts of residential-edge traversals) into one scalar; a route can
appear optimal merely because the chosen weights hide the safety
cost. Second, treating risk as exogenous and uncalibrated means a
practitioner cannot interpret the resulting plan as ``safe'' in any
defensible operational sense.

\section{Problem Statement}

This thesis develops an SA-VRPTW formulation in which service quality
is the primary optimisation target, while crash-survival risk and
residential exposure are enforced as explicit, auditable
$\varepsilon$-constraints. The associated software pipeline materialises
realistic instances from public geographic and road-safety data, and
the solver stack reports $F_1$ through a single evaluation contract
shared across exact and heuristic methods.

\section{Contributions}

The contributions of this thesis are:

\begin{enumerate}
    \item A Safety-Aware VRPTW formulation that treats per-route
    crash exposure as a hard $\varepsilon$-constraint
    (survival probability $\geq e^{-\bar{R}}$) rather than as an
    objective component, so the safety/efficiency trade-off is traced
    \emph{explicitly} as a Pareto frontier
    \citep{salmon2023,ehrgott2005} rather than implicitly through a
    weighted sum.
    \item A fleet-wide budget on residential-street penetration plus
    a per-route residential cap, so a planner can constrain
    network-level use of vulnerable streets even when individual
    routes are safe.
    \item A documented public-source crash-risk calibration chain
    based on \citet{morth2022} category-level fatality totals,
    \citet{mohan2017} ordering constraints, and an OpenStreetMap
    structural proxy, avoiding the protected Integrated Road Accident
    Database microdata.
    \item A convex SLA penalty $\exp(\beta_{\text{stw}}\tau)-1$
    replacing the linear lateness term common in soft-time-window
    formulations \citep{koskosidis1992,fu2008,taillard1997},
    approximated via piecewise-linear outer envelope in the MILP and
    evaluated exactly in the heuristics.
    \item A unified solver/evaluation contract under which an exact
    MILP, a giant-tour genetic algorithm with Bellman--Ford optimal
    split decoding \citep{prins2004}, and an Adaptive Large
    Neighbourhood Search \citep{ropke2006} all evaluate solutions
    through identical $F_1$ and feasibility code, eliminating the
    apples-to-pears bias common in VRP solver comparisons.
    \item A five-city empirical protocol on real Blinkit dark-store
    coordinates extracted from a publicly committed KML snapshot,
    rather than synthetic uniform-disk depots, so that the geometric
    difficulty of the instances reflects the actual topology
    operators face.
\end{enumerate}

\section{Thesis Organisation}

The remainder of this thesis is organised as follows. Chapter~2
reviews related work on VRPTW with soft time windows, safety-aware
routing, multi-objective $\varepsilon$-constraint methods,
quick-commerce dispatch, and crash-risk calibration. Chapter~3 states
the SA-VRPTW formulation in full, including the multi-depot coupling
and the MILP linearisation of the convex SLA penalty. Chapter~4
documents the calibration and data pipeline---the BASM crash-risk
chain, the BPR congestion model, the dark-store snapshot, and the
OpenStreetMap cross-validation. Chapter~5 details the three solvers
together with formal pseudocode. Chapter~6 defines the experimental
protocol, the Pareto sweep, and the post-hoc Monte Carlo evaluators.
Chapter~7 concludes with operator-facing implications, limitations,
and directions for future work.
